% ID: SEL-D0049-OB-E04
% ESTADO: cotejada. Segunda pasada visual y borrador compilado revisados.
% Fuente: Selectividad/sel_2013_matematicas_aplicadas/junio/reserva b/EXAMEN 1-ANDALUCÍA.pdf
% SHA-256: ec6fd00d0daee547aa6128f3f6e1099278ee551f7f80917b331a1772764317b3
% Página: 2; opción B; ejercicio original 4.

\begin{enumerate}[label=\alph*)]
\item \textbf{(1.25 puntos)} Una población de 6000 personas se ha dividido en 3 estratos, uno con 1000 personas, otro con 3500 y otro con 1500. En esa población se ha realizado un muestreo estratificado con afijación proporcional, en el que se han elegido al azar 15 personas del tercer estrato. Determine el tamaño de la muestra total obtenida con este muestreo y su composición.
\item \textbf{(1.25 puntos)} Dada la población $\{1,4,7\}$, construya todas las muestras posibles de tamaño 2 que puedan formarse mediante muestreo aleatorio simple, y halle la varianza de las medias muestrales de todas esas muestras.
\end{enumerate}
